\documentclass[11pt]{article}
\usepackage[margin=1in]{geometry}
\usepackage{amsmath,amssymb,amsthm}
\usepackage{hyperref}
\usepackage{enumitem}
\usepackage{xurl}
\let\oldus\_
\renewcommand{\_}{\oldus\allowbreak}

\newtheorem{theorem}{Theorem}
\newtheorem{lemma}[theorem]{Lemma}
\newtheorem{corollary}[theorem]{Corollary}
\theoremstyle{definition}
\newtheorem{definition}[theorem]{Definition}
\theoremstyle{remark}
\newtheorem{remark}[theorem]{Remark}

\title{A disproof of Conjecture 00000004053\\[2pt]
\large The trivial $\mathbb{F}_3[C_3]$-module has syzygy period $2$, and $2$ divides no prime factor of the order of any subgroup of $C_3$}
\date{}

\begin{document}
\sloppy
\maketitle

\section{The conjecture}

Conjecture 00000004053 reads (English):
\begin{quote}
\emph{Definition:} A periodic module is one with $\Omega^n M$ isomorphic to $M$.
\emph{Conjecture:} The prime factors of the minimal period of a periodic module always divide
prime factors of the order of the defect group, and every power of every admissible prime is
realized in some explicit block. (period-defect group prime factor correspondence)
\end{quote}
The Chinese text says the same: the prime factors of the minimal period of a periodic module
always divide prime factors of the order of the defect group, and every power of every admissible
prime is realized in some explicit block.

\textbf{Reading.} The setting is modular representation theory: a finite group $G$, a field $k$ of
characteristic $p$, a block of $kG$ with its defect group $D$, and a periodic module $M$ in that
block, i.e.\ a finitely generated module with $\Omega^n M\cong M$ for some $n\ge 1$, where
$\Omega M$ is the kernel of a projective cover of $M$ (the Heller operator). The \emph{minimal
period} is the least such $n$. Clause 1 says: every prime factor $q$ of the minimal period divides
some prime factor $r$ of $|D|$. Since $q$ and $r$ are primes, this is the same as $q=r$, and the
same as $q\mid |D|$; we refute both formulations. Clause 2 (``admissible prime'', ``explicit
block'') is not defined in the text and is not used: the conjecture is a conjunction, and we
refute Clause 1.

\textbf{Counterexample.} $p=3$, $k=\mathbb{F}_3$, $G=C_3$, $M=k$ the trivial module. Its minimal
period is $2$, while every subgroup of $C_3$ (in particular every possible defect group) has order
$1$ or $3$; $2$ divides neither $3$ nor any prime factor of $1$ or $3$.

\section{Definitions (as in the Lean file)}

All modules are left modules over a ring $R$; in Lean, modules are objects of
\texttt{ModuleCat.\{0\} R}.
\begin{itemize}[leftmargin=1.5em]
\item A submodule $K\le P$ is \emph{superfluous} if $K+L=P$ implies $L=P$ (\texttt{IsSuperfluous}).
\item $\pi:P\to M$ is a \emph{projective cover} if $P$ is projective, $\pi$ is onto, and
$\ker\pi$ is superfluous (\texttt{IsProjectiveCover}). This is the definition retrieved from
Wikipedia (\emph{Projective cover}): an epimorphism $p:P\to M$ ``such that the kernel of p is a
superfluous submodule of P.''
\item $N$ is a \emph{syzygy} of $M$ (\texttt{IsSyzygy R M N}) if $N\cong\ker\pi$ for some projective
cover $\pi:P\to M$. \texttt{IterSyz R n M N} means $N\cong\Omega^n M$: there is a chain
$M=M_0,M_1,\dots,M_n$ with $M_{i+1}$ a syzygy of $M_i$ and $M_n\cong N$.
\item \texttt{HasPeriod R M n}: $n\ge1$, $\Omega^n M\cong M$, and $\Omega^m M\not\cong M$ for
$0<m<n$ (for every chain). So $n$ is the minimal period.
\item $M$ is \emph{indecomposable} (\texttt{IsIndecomposable}) if $M\ne0$ and whenever
$M=S\oplus T$ for submodules $S,T$ (\texttt{IsCompl S T}), $S=0$ or $T=0$.
\item Defect group. Wikipedia (\emph{Modular representation theory}): ``To each block B of the group
algebra K[G], Brauer associated a certain p-subgroup, known as its defect group (where p is the
characteristic of K).'' We use only that a defect group is a subgroup of $G$. (Wikipedia also states
that ``the Sylow p-subgroup of the finite group G is a defect group for the principal block of
K[G]''; for $G=C_3$, $p=3$ this is $C_3$, but we do not need it.)
\end{itemize}

\textbf{The group algebra.} $G=C_3$ is Mathlib's \texttt{Multiplicative (ZMod 3)} with generator
$g$; $A=\mathbb{F}_3[C_3]$ is Mathlib's group algebra \texttt{MonoidAlgebra (ZMod 3) G}. Let
$t=g-1\in A$ and let $\varepsilon:A\to\mathbb{F}_3$ be the augmentation ($\varepsilon(h)=1$ for
$h\in G$, Mathlib's \texttt{MonoidAlgebra.lift} of the trivial character). The trivial module is
$k=\mathbb{F}_3$ with $a\cdot x=\varepsilon(a)x$ (\texttt{Triv}); every $h\in G$ acts as the identity
(\texttt{triv\_group\_acts\_trivially}).

\section{Proof}

\begin{lemma}[\texttt{t\_cube}, \texttt{t\_sq\_ne\_zero}, \texttt{exists\_t\_mul}]
$t^3=0$, $t^2\ne0$, and every $a\in A$ is $a=\varepsilon(a)+td$ for some $d\in A$; so
$\ker\varepsilon=tA$.
\end{lemma}
\begin{proof}
$(g-1)^3=g^3-1-3(g^2-g)=0$ since $g^3=1$ and $3=0$ in $A$. Also $t^2=g^2-2g+1=g^2+g+1$, whose
coefficient at $1\in G$ is $1$. For the decomposition it suffices (by linearity) to treat $a=h\in G$:
$h=g^m$ and $g^m-1=(g-1)(1+g+\dots+g^{m-1})$.
\end{proof}

\begin{remark}
Hence $1,t,t^2$ span $A$, and since $\dim A=3$ they form a basis; with $t^3=0$ this gives
$A\cong\mathbb{F}_3[t]/(t^3)$ via $t=g-1$. The Lean proof does not use this isomorphism; it works in
the group algebra itself, using only the three facts of Lemma 1.
\end{remark}

\begin{lemma}\label{lem:alg}
(a) $t^2c=\varepsilon(c)t^2$ for all $c$ (\texttt{t\_sq\_mul}); hence $t^2c=0\Rightarrow\varepsilon(c)=0$.
(b) If $a\ne0$ then $t^2\in Aa$ (\texttt{t\_sq\_mem}).
(c) Every submodule of $A$ inside $\ker\varepsilon$ is superfluous (\texttt{superfluous\_of\_le\_ker}).
(d) If $P$ is projective and $t^2p=0$ then $p\in tP$ (\texttt{mem\_t\_of\_t\_sq\_smul}).
\end{lemma}
\begin{proof}
(a) $c=\varepsilon(c)+td$ and $t^3=0$. (b) If $\varepsilon(a)\ne0$, $t^2=\varepsilon(a)^{-1}t^2a$.
Otherwise $a=td$; if $\varepsilon(d)\ne0$ then $ta=\varepsilon(d)t^2$; otherwise $a=t^2d'=\varepsilon(d')t^2$
with $\varepsilon(d')\ne0$. (c) If $K+L=A$ then $1=x+l$ with $\varepsilon(l)=1$, $l=1+td$, and
$(1-td+t^2d^2)l=1$, so $L=A$. (d) $P$ is a direct summand of a free module
(\texttt{Module.projective\_def'}); the coordinates $c_i$ of $p$ satisfy $t^2c_i=0$, so
$c_i\in tA$ by (a) and Lemma 1, and $p\in tP$.
\end{proof}

\begin{lemma}[cyclic modules]\label{lem:cyc}
Let $M=Am_0$ with $m_0\notin tM$ (\texttt{CyclicGen}). Then
(i) $A\to M$, $a\mapsto am_0$, is a projective cover (\texttt{toSpan\_isCover});
(ii) for \emph{every} projective cover $\pi:P\to M$, $\ker\pi\cong\operatorname{ann}(m_0)$
(\texttt{ker\_cover\_equiv}).
\end{lemma}
\begin{proof}
(i) If $am_0=0$ and $\varepsilon(a)\ne0$, writing $a=\varepsilon(a)+td$ gives
$m_0=t\cdot(-\varepsilon(a)^{-1}d\,m_0)\in tM$; so the kernel lies in $\ker\varepsilon$ and is
superfluous by Lemma \ref{lem:alg}(c). (ii) Pick $p_0$ with $\pi(p_0)=m_0$ and let
$\varphi(a)=ap_0$. Then $\operatorname{im}\varphi+\ker\pi=P$, so $\varphi$ is onto
(superfluous kernel). If $ap_0=0$ with $a\ne0$, then $t^2p_0=0$ by Lemma \ref{lem:alg}(b), so
$p_0=tq$ by (d) and $m_0=t\pi(q)$, a contradiction; so $\varphi$ is injective. Thus $\varphi:A\cong P$,
$\pi\circ\varphi$ is $a\mapsto am_0$, and $\varphi$ maps $\operatorname{ann}(m_0)$ onto $\ker\pi$.
\end{proof}

\begin{theorem}[\texttt{period\_K}, \texttt{iterSyz\_two\_any}]
The trivial module $k$ is finitely generated, simple (hence indecomposable) and not projective; its
minimal syzygy period is $2$; and every second syzygy of $k$, for every choice of projective covers,
is isomorphic to $k$.
\end{theorem}
\begin{proof}
$k=A\cdot1$ and $1\notin tk=0$. By Lemma \ref{lem:cyc}, $\Omega k\cong N_1=\operatorname{ann}(1)=\ker\varepsilon$
for every cover. $N_1=A\cdot t$ and $t\notin tN_1$ (else $t=t^2d$ and $t^2=t^3d=0$), so every
syzygy of $N_1$ (or of any module isomorphic to it, \texttt{CyclicGen.map}) is isomorphic to
$N_2=\operatorname{ann}(t)=\{a: at=0\}$. The map $x\mapsto xt^2$ is an $A$-isomorphism $k\to N_2$
(\texttt{trivEquivN2}): it is $A$-linear by Lemma \ref{lem:alg}(a), injective since $t^2\ne0$, and onto
since $at=0$ forces $a=t^2d'=\varepsilon(d')t^2$. Finally $N_1\not\cong k$ since $t\cdot t=t^2\ne0$ in
$N_1$ while $t$ kills $k$ (\texttt{N1\_not\_equiv\_triv}); so no first syzygy of $k$ is isomorphic to
$k$ (\texttt{not\_iterSyz\_one}). Simplicity: any nonzero $x\in k$ generates $k$ (\texttt{triv\_simple}).
Non-projectivity: $t^2\cdot1=0$ but $1\notin tk$, contradicting Lemma \ref{lem:alg}(d)
(\texttt{triv\_not\_projective}).
\end{proof}

\begin{theorem}[\texttt{conjecture\_4053\_false}, \texttt{conjecture\_4053\_false'}]
Let \texttt{Clause1 p H} say: for every finitely generated, indecomposable, non-projective
$\mathbb{F}_p[H]$-module $M$ with minimal period $n$ there is a subgroup $D\le H$ such that every prime
factor $q$ of $n$ divides some prime factor $r$ of $|D|$ (\texttt{Clause1'}: $q\mid|D|$). Then
\texttt{Clause1 3 C\_3} and \texttt{Clause1' 3 C\_3} are false.
\end{theorem}
\begin{proof}
Take $M=k$, $n=2$, $q=2$. By Lagrange (\texttt{Subgroup.card\_subgroup\_dvd\_card}), $|D|\mid 3$ for every
$D\le C_3$, so $2\nmid|D|$ and $2$ divides no prime factor of $|D|$ (\texttt{not\_two\_dvd\_card}).
\end{proof}

Since the defect group of any block of $\mathbb{F}_3[C_3]$ is a subgroup of $C_3$, Clause 1 of the
conjecture implies \texttt{Clause1 3 C\_3}; hence the conjecture is false
(\texttt{conjecture\_4053\_conj\_false}: $\neg(\texttt{Clause1 3 C\_3}\wedge Q)$ for every $Q$).

\section{Relation to the earlier PR \#291}

PR \#291 used the same counterexample ($\mathbb{F}_3[C_3]$, trivial module, period $2$, defect
group of order $3$) and was closed without merging because its Lean file was vacuous: its main
theorem was the arithmetic statement $1=1\wedge 2=1+1\wedge 3=3\wedge 3\bmod 2=1\wedge 4=2\cdot2$, and
the review noted that no module, group algebra, $\Omega$-operator or defect group existed in Lean.
The present formalization contains the actual objects: Mathlib's group algebra
$\mathbb{F}_3[C_3]$, the trivial module with the group acting trivially, projective covers with
superfluous kernels, syzygies as kernels of projective covers (for arbitrary projective modules, not a
fixed choice), the computation $\Omega k\cong\ker\varepsilon$, $\Omega^2k\cong k$, $\Omega k\not\cong k$,
the minimal-period predicate, and the quantifier over all subgroups of $C_3$.

\section{Scope}

Refuted: Clause 1 (in both formulations above) for $p=3$, $G=C_3$, $k=\mathbb{F}_3$, with the
defect group replaced by an arbitrary subgroup of $C_3$, and with periodic modules restricted to
finitely generated indecomposable non-projective modules; hence the conjunction. Conventions:
syzygies are a relation ($N\cong\ker\pi$ for some projective cover $\pi$); we prove that the period
does not depend on the choice of covers. Modules in Lean live in \texttt{Type} (universe 0); the
projective cover $P$ in a syzygy may be any projective module, not only a finitely generated one.
Not formalized: blocks and defect groups themselves (only the fact that a defect group is a
subgroup of $G$ is used), Clause 2, and the reading in which $\Omega^nM\cong M$ is taken in the stable
module category. The same argument works for every odd prime $p$ with $G=C_p$ (not formalized).

\section{Lean formalization}

File \texttt{lean/Conjecture4053/Basic.lean} (namespace \texttt{C4053}), only import \texttt{Mathlib},
Lean toolchain and Mathlib as in \texttt{lean-toolchain}/\texttt{lake-manifest.json}. Main statements:
\begin{itemize}[leftmargin=1.5em]
\item \texttt{conjecture\_4053\_false : \textasciitilde Clause1 3 G} and
\texttt{conjecture\_4053\_false' : \textasciitilde Clause1' 3 G};
\item \texttt{trivial\_module\_counterexample}: \texttt{Module.Finite A Triv},
\texttt{IsIndecomposable A Triv}, \texttt{\textasciitilde Module.Projective A Triv},
\texttt{HasPeriod A K 2}, every \texttt{IterSyz A 2 K N} has \texttt{N} isomorphic to \texttt{Triv}, and
for every \texttt{D : Subgroup G} and every prime factor \texttt{r} of \texttt{Nat.card D},
\texttt{\textasciitilde 2 | r};
\item \texttt{conjecture\_4053\_conj\_false (Clause2 : Prop) : \textasciitilde(Clause1 3 G /\textbackslash{} Clause2)}.
\end{itemize}
Axiom audit (\texttt{\#print axioms} for all four): \texttt{[propext, Classical.choice, Quot.sound]}.
No \texttt{sorry}, no \texttt{native\_decide}, no added axioms. Built with \texttt{lake build}.

\section*{Sources}
\begin{itemize}[leftmargin=1.5em]
\item Wikipedia, \emph{Projective cover}, \url{https://en.wikipedia.org/wiki/Projective_cover}:
``such that the kernel of p is a superfluous submodule of P.''
\item Wikipedia, \emph{Modular representation theory},
\url{https://en.wikipedia.org/wiki/Modular_representation_theory}: ``To each block B of the group
algebra K[G], Brauer associated a certain p-subgroup, known as its defect group (where p is the
characteristic of K).''; ``the Sylow p-subgroup of the finite group G is a defect group for the
principal block of K[G].''
\item GitHub PR \#291 (closed), review comment:
\url{https://github.com/The-Last-Math-Competition/The-Last-Math-Competition/pull/291}.
\end{itemize}

\end{document}
